% ═══════════════════════════════════════════════════════════════
% MUSTERLÖSUNG: Las direcciones (Wegbeschreibung)
% ═══════════════════════════════════════════════════════════════
% Sequenz 6c: Mi ciudad - Las direcciones
% Fachfarbe: #c41e3a (Karminrot)
% Lösungsfarbe: ROT
% ═══════════════════════════════════════════════════════════════

\documentclass[11pt, a4paper]{article}

% ═══════════════════════════════════════════════════════════════
% SW-MODUS TOGGLE
% ═══════════════════════════════════════════════════════════════
\newif\ifswmodus
\swmodusfalse

% ═══════════════════════════════════════════════════════════════
% PAKETE
% ═══════════════════════════════════════════════════════════════

\usepackage[utf8]{inputenc}
\usepackage[T1]{fontenc}
\usepackage[ngerman]{babel}
\usepackage{helvet}
\renewcommand{\familydefault}{\sfdefault}

\usepackage[margin=2cm]{geometry}
\usepackage{enumitem}
\usepackage{tabularx}
\usepackage{booktabs}
\usepackage{multirow}

\usepackage{xcolor}
\usepackage{framed}
\usepackage{tcolorbox}
\tcbuselibrary{skins}

\usepackage{fancyhdr}
\usepackage{lastpage}
\usepackage{needspace}
\usepackage{tikz}

% ═══════════════════════════════════════════════════════════════
% FARBEN
% ═══════════════════════════════════════════════════════════════

\definecolor{spanisch-color}{HTML}{c41e3a}
\definecolor{grau-color}{HTML}{888888}
\definecolor{hellgrau-color}{HTML}{f5f5f5}
\definecolor{orange-color}{HTML}{b35c00}
\definecolor{loesung-color}{HTML}{c62828}

\definecolor{sw-dark}{HTML}{000000}
\definecolor{sw-gray}{HTML}{666666}
\definecolor{sw-white}{HTML}{ffffff}

\ifswmodus
    \colorlet{spanisch}{sw-dark}
    \colorlet{grau}{sw-gray}
    \colorlet{hellgrau}{sw-white}
    \colorlet{orange}{sw-dark}
    \colorlet{loesung}{sw-dark}
\else
    \colorlet{spanisch}{spanisch-color}
    \colorlet{grau}{grau-color}
    \colorlet{hellgrau}{hellgrau-color}
    \colorlet{orange}{orange-color}
    \colorlet{loesung}{loesung-color}
\fi

\newcommand{\fachfarbe}{spanisch}

% ═══════════════════════════════════════════════════════════════
% VARIABLEN
% ═══════════════════════════════════════════════════════════════

\newcommand{\thema}{Las direcciones -- Wegbeschreibung}
\newcommand{\sequenz}{06}
\newcommand{\block}{c}
\newcommand{\fach}{Spanisch}
\newcommand{\klasse}{7}
\newcommand{\maxpunkte}{20}
\newcommand{\datum}{\today}

% ═══════════════════════════════════════════════════════════════
% KOPF- UND FUSSZEILE
% ═══════════════════════════════════════════════════════════════

\pagestyle{fancy}
\fancyhf{}
\renewcommand{\headrulewidth}{0.4pt}
\renewcommand{\footrulewidth}{0.4pt}

\lhead{\fach{} -- Klasse \klasse}
\chead{\textcolor{loesung}{\textbf{ML-\sequenz\block{} (MUSTERLÖSUNG)}}}
\rhead{\datum}

\lfoot{}
\cfoot{}
\rfoot{Seite \thepage\ von \pageref{LastPage}}

% ═══════════════════════════════════════════════════════════════
% BEFEHLE
% ═══════════════════════════════════════════════════════════════

\newcommand{\punktebox}[1]{\hfill\fbox{\small \_/#1}}
\newcommand{\loesung}[1]{\textcolor{loesung}{\textbf{#1}}}
\newcommand{\luecke}[1]{\rule{#1}{0.4pt}}

\newtcolorbox{merkkasten}[1][]{
    colback=hellgrau,
    colframe=\fachfarbe,
    colbacktitle=white,
    coltitle=\fachfarbe,
    boxrule=1pt,
    arc=0pt,
    left=8pt, right=8pt, top=8pt, bottom=8pt,
    fonttitle=\bfseries,
    attach boxed title to top left={yshift=-2mm, xshift=5mm},
    boxed title style={boxrule=0pt, colback=white},
    title=#1
}

\newtcolorbox{vokabelkasten}{
    colback=white,
    colframe=\fachfarbe,
    boxrule=0.5pt,
    arc=0pt,
    left=5pt, right=5pt, top=5pt, bottom=5pt
}

\newtcolorbox{hinweis}{
    colback=white,
    colframe=orange,
    colbacktitle=white,
    coltitle=orange,
    boxrule=1pt,
    arc=0pt,
    left=5pt, right=5pt, top=5pt, bottom=5pt,
    fonttitle=\bfseries,
    attach boxed title to top left={yshift=-2mm, xshift=5mm},
    boxed title style={boxrule=0pt, colback=white},
    title=Hinweis
}

\newenvironment{aufgabe}[1]{%
    \needspace{5\baselineskip}%
    \nopagebreak[4]%
    \vspace{10pt}
    \noindent\textbf{\textcolor{\fachfarbe}{Aufgabe #1}}
    \nopagebreak[4]%
    \vspace{5pt}
    \nopagebreak[4]%
    \begin{list}{}{\leftmargin=0pt}
    \item[]
}{%
    \end{list}
}

\newcommand{\es}[1]{\textcolor{\fachfarbe}{\textbf{#1}}}

% ═══════════════════════════════════════════════════════════════
% DOKUMENT
% ═══════════════════════════════════════════════════════════════

\begin{document}

% ═══════════════════════════════════════════════════════════════
% KOPFBEREICH
% ═══════════════════════════════════════════════════════════════

\begin{center}
    {\Large\bfseries\textcolor{\fachfarbe}{Arbeitsblatt: \thema}}\\[0.3em]
    {\large\textcolor{loesung}{\textbf{--- MUSTERLÖSUNG ---}}}
\end{center}

\vspace{5pt}

\noindent
\begin{tabularx}{\textwidth}{|X|X|X|}
\hline
\textbf{Name:} \textcolor{loesung}{Musterschüler/in} &
\textbf{Klasse:} \klasse &
\textbf{Datum:} \datum \\
\hline
\end{tabularx}

\vspace{15pt}

% ═══════════════════════════════════════════════════════════════
% MERKKASTEN: Wegbeschreibung (mit Lösungen)
% ═══════════════════════════════════════════════════════════════

\begin{merkkasten}[Merkkasten: Wegbeschreibung]
\textbf{Richtungsangaben (las direcciones):}

\vspace{0.5em}

\begin{tabular}{ll}
\es{todo recto} & = \loesung{geradeaus} \\[0.8em]
\es{a la derecha} & = \loesung{nach rechts} \\[0.8em]
\es{a la izquierda} & = \loesung{nach links} \\[0.8em]
\es{girar} & = \loesung{abbiegen} \\
\end{tabular}

\vspace{1em}

\textbf{Ortsangaben (ubicación):}

\vspace{0.5em}

\begin{tabular}{ll}
\es{al lado de} & = \loesung{neben} \\[0.8em]
\es{enfrente de} & = \loesung{gegenüber von} \\[0.8em]
\es{cerca de} & = \loesung{in der Nähe von} \\[0.8em]
\es{lejos de} & = \loesung{weit weg von} \\
\end{tabular}

\vspace{1em}

\textbf{Fragen:}

\vspace{0.5em}

\begin{tabular}{ll}
\es{¿Dónde está...?} & = \loesung{Wo ist...?} \\[0.8em]
\es{¿Cómo llego a...?} & = \loesung{Wie komme ich zu...?} \\
\end{tabular}

\end{merkkasten}

\vspace{15pt}

% ═══════════════════════════════════════════════════════════════
% AUFGABE 1: Zuordnung (mit Lösungen)
% ═══════════════════════════════════════════════════════════════

\begin{aufgabe}{1}
Verbinde Spanisch und Deutsch. Schreibe die richtige Nummer in die Klammer. \punktebox{4}
\end{aufgabe}

\begin{center}
\begin{tabular}{rlcrl}
1. & \es{todo recto} & \hspace{2cm} & ( \loesung{3} ) & nach links \\[0.6em]
2. & \es{a la derecha} & & ( \loesung{4} ) & in der Nähe von \\[0.6em]
3. & \es{a la izquierda} & & ( \loesung{1} ) & geradeaus \\[0.6em]
4. & \es{cerca de} & & ( \loesung{2} ) & nach rechts \\
\end{tabular}
\end{center}

\vspace{15pt}

% ═══════════════════════════════════════════════════════════════
% AUFGABE 2: Lückentext (mit Lösungen)
% ═══════════════════════════════════════════════════════════════

\begin{aufgabe}{2}
Ergänze die Sätze mit den passenden Wörtern. \punktebox{4}
\end{aufgabe}

\begin{vokabelkasten}
\textit{Wortliste:} enfrente de | al lado de | lejos de | todo recto
\end{vokabelkasten}

\vspace{0.5em}

\noindent
a) La farmacia está \es{\loesung{enfrente de}} la panadería. (= gegenüber) \\[0.8em]
b) El banco está \es{\loesung{al lado de}} el supermercado. (= neben) \\[0.8em]
c) La estación está \es{\loesung{lejos de}} de aquí. (= weit weg) \\[0.8em]
d) Sigue \es{\loesung{todo recto}} y llegas a la plaza. (= geradeaus)

\vspace{15pt}

% ═══════════════════════════════════════════════════════════════
% AUFGABE 3: Dialog vervollständigen (mit Lösungen)
% ═══════════════════════════════════════════════════════════════

\begin{aufgabe}{3}
Ein Tourist fragt nach dem Weg. Vervollständige den Dialog. \punktebox{6}
\end{aufgabe}

\noindent
\textbf{Turista:} \es{Perdone, ¿cómo llego a la estación?}

\vspace{0.3em}

\textbf{Tú:} \loesung{Sigue todo recto. / Sigue recto por esta calle.}

\vspace{0.8em}

\textbf{Turista:} \es{¿Y después?}

\vspace{0.3em}

\textbf{Tú:} \loesung{Gira a la derecha/izquierda. / La estación está al lado del supermercado.}

\vspace{0.8em}

\textbf{Turista:} \es{¿Está lejos?}

\vspace{0.3em}

\textbf{Tú:} \loesung{No, está cerca. / Sí, está un poco lejos.}

\vspace{0.8em}

\textit{\textcolor{grau}{Bewertung: 2P pro sinnvolle Antwort mit korrekter Struktur}}

\vspace{15pt}

% ═══════════════════════════════════════════════════════════════
% AUFGABE 4: Wegbeschreibung mit Stadtplan (mit Lösungen)
% ═══════════════════════════════════════════════════════════════

\begin{aufgabe}{4}
Schau dir den Stadtplan an. Beschreibe den Weg vom Punkt \textbf{A} (Start) zur \textbf{Farmacia}. Schreibe mindestens 3 Sätze. \punktebox{6}
\end{aufgabe}

\begin{center}
\begin{tikzpicture}[scale=0.8]
% Straßen
\draw[thick] (0,0) -- (10,0);
\draw[thick] (0,3) -- (10,3);
\draw[thick] (0,6) -- (10,6);
\draw[thick] (2,0) -- (2,6);
\draw[thick] (5,0) -- (5,6);
\draw[thick] (8,0) -- (8,6);

% Straßennamen
\node[above] at (5,6) {\small Calle Mayor};
\node[above] at (5,3) {\small Calle del Sol};
\node[above] at (5,0) {\small Calle de la Luna};

% Gebäude
\node[draw, fill=white, minimum width=1.5cm, minimum height=1cm] at (0.8,4.5) {\tiny Banco};
\node[draw, fill=white, minimum width=1.5cm, minimum height=1cm] at (3.5,4.5) {\tiny Panadería};
\node[draw, fill=white, minimum width=1.5cm, minimum height=1cm] at (6.5,4.5) {\tiny Farmacia};
\node[draw, fill=white, minimum width=1.5cm, minimum height=1cm] at (9.2,4.5) {\tiny Hotel};

\node[draw, fill=white, minimum width=1.5cm, minimum height=1cm] at (0.8,1.5) {\tiny Super-};
\node[draw, fill=white, minimum width=1.5cm, minimum height=1cm] at (0.8,1) {\tiny mercado};
\node[draw, fill=white, minimum width=1.5cm, minimum height=1cm] at (3.5,1.5) {\tiny Estación};
\node[draw, fill=white, minimum width=1.5cm, minimum height=1cm] at (6.5,1.5) {\tiny Correos};
\node[draw, fill=white, minimum width=1.5cm, minimum height=1cm] at (9.2,1.5) {\tiny Biblioteca};

% Startpunkt
\node[circle, fill=spanisch-color, text=white, minimum size=0.6cm] at (0.5,0) {\textbf{A}};

% Legende
\node at (5,-1) {\small A = Startpunkt (Tú estás aquí)};

% Lösungspfad einzeichnen
\draw[->, very thick, loesung-color, dashed] (0.5,0.3) -- (0.5,3) -- (0.5,6) -- (6.5,6) -- (6.5,5);
\end{tikzpicture}
\end{center}

\vspace{0.5em}

\textbf{Beispiellösung:}

\loesung{Sigue todo recto por la Calle de la Luna.}

\loesung{Gira a la derecha en la Calle Mayor.}

\loesung{La farmacia está a la derecha, al lado de la panadería.}

\loesung{(Alternativ: La farmacia está enfrente de Correos.)}

\vspace{0.5em}

\textit{\textcolor{grau}{Bewertung: 2P pro korrekten Satz (Richtung + Struktur)}}

% ═══════════════════════════════════════════════════════════════
% PUNKTEÜBERSICHT
% ═══════════════════════════════════════════════════════════════

\vspace{20pt}
\hrule
\vspace{10pt}

\begin{flushright}
\textbf{Punkte:} \loesung{\maxpunkte} / \maxpunkte
\end{flushright}

\end{document}
